\chapter{Applications}\label{ch:iv:applications}

Four hundred CVs arrive for twenty first-round slots, and the reader has one afternoon: thirty-six seconds a
CV. Whatever is not on the first half of the first page is, for practical purposes, not there. A line that
says what was built and what number it moved is read; a line that lists technologies is skimmed; a line that
claims a Sharpe ratio of 3 on a student project is read twice, for the wrong reason. An application is the
first measurement in the funnel of \cref{ch:iv:what-each-interview-tests-by-role}, and the cheapest one to
improve.

\section{The one-page CV}

\begin{definition}[CV screen]\label{def:iv:applications:screen}
The \emph{CV screen}\index{CV screen} is the stage at which a person or a program reads an application and
decides whether it goes on to an assessment or an interview. Its inputs are the CV, any form answers and,
where the firm has them, a referral and the candidate's earlier applications.
\end{definition}

The reader of a screen has a short list of questions: is the candidate eligible (degree, date, right to
work), is there evidence of the role's core abilities, and is there anything exceptional. The first half-page
has to answer all three.

\begin{method}[Writing a CV line]\label{met:iv:applications:line}
\begin{enumerate}
\item Start with a verb for what you did (built, measured, proved, reduced), not with the tools.
\item Name the object: a pricer, a feed parser, a study of an effect.
\item Give a number with its baseline: ``from 40 to 6 milliseconds'', ``out-of-sample $R^2$ of 0.8\% against
0.1\% for the benchmark'', ``1\,200 tests''.
\item Name the tools last, and only those you can be interviewed on for thirty minutes.
\item Cut every line that fails steps 1 to 3. One page; the most relevant evidence first.
\end{enumerate}
\end{method}

\begin{example}[One line, before and after]\label{ex:iv:applications:line}
Before: ``Used Python, pandas and scikit-learn to build machine-learning models for stock prediction.''
After: ``Tested whether overnight order imbalance predicts next-day returns in 500 US stocks, 2015--2024;
walk-forward $R^2$ 0.3\% ($t = 2.1$), gone after costs; Python.'' The second line tells a research
interviewer what to ask about, and its honest last clause is worth more than the result: it shows the
candidate knows what costs and multiple testing do (\cref{ch:iv:statistics}).
\end{example}

A CV is also the first list of interview questions. Every line invites ``tell me about this'', and the
answer should be ready in two lengths: ten seconds, and five minutes with the numbers and what went wrong.

\section{Internships, projects and competitions}

Evidence of the role's abilities comes in three forms, each read differently. An internship at a trading firm
or bank is evidence that another firm's screen was passed and that the candidate worked in the environment;
its value is in what was delivered, which the candidate must be able to describe without disclosing the
firm's confidential information. A personal project is evidence of initiative, and its value is in the
honesty of its measurement: a baseline, an out-of-sample test, a statement of what failed. A competition
(mathematics olympiads, programming contests, data-science and trading competitions; the pipelines are in
One Quant Book~17, chapter~29) is evidence of ability under time pressure, read through its ranking.

Rankings in data competitions need care, because a public leaderboard scored on a small sample rewards noise.

\begin{proposition}[The best of $k$ equal scores]\label{prop:iv:applications:max}
If $k$ teams of equal true skill are scored with independent noise of standard deviation $\sigma$, the
expected best score exceeds the common true score by $\sigma\,m_k$, where $m_k = \E[\max_{i\le k} Z_i]$ for
independent standard normal $Z_i$: $m_{10} \approx 1.54$, $m_{60} \approx 2.32$, $m_{100} \approx 2.51$,
$m_{400} \approx 2.97$. The asymptotic $\sqrt{2\ln k}$ overstates at these sizes ($3.46$ for $k = 400$).
\end{proposition}

\begin{proof}
The maximum has density $k\,\varphi(x)\,\Phi(x)^{k-1}$; the values are its mean by numerical integration,
confirmed by simulation in the chapter's tests.
\end{proof}

The same proposition explains why a strong result from many attempts is weaker evidence than it looks: the
best of sixty backtests of worthless variants over five years has an expected annualised Sharpe ratio of about
$2.32/\sqrt{5} \approx 1.04$ (One Quant Book~4, chapter~12, on multiple testing and the deflated Sharpe
ratio). A candidate who reports how many variants were tried is believed more, not less.

\section{Referrals and recruiters}

\begin{definition}[Employee referral, agency recruiter]\label{def:iv:applications:referral}
An \emph{employee referral}\index{employee referral} is an application submitted or endorsed by a current
employee of the firm, who vouches that the candidate is worth screening. An \emph{agency
recruiter}\index{agency recruiter} is an intermediary retained by the hiring firm, usually paid a fee on a
successful hire, who finds candidates and manages their process on the firm's behalf.
\end{definition}

A referral changes the screen's prior. In a study of nine large firms in three industries, referred applicants
were more likely to be hired and to accept offers, with similar skills on paper, and referred workers were
substantially less likely to quit (Burks, Cowgill, Hoffman and Housman, 2015). The mechanism is information:
the referrer knows the candidate and the firm. A referral from someone who has not worked with the candidate
carries little of it, and a good referrer says so.

A recruiter is paid by the firm, so the candidate is the product being placed, not the client. That does not
make the recruiter an adversary: a placement requires the candidate to accept, and a recruiter who places
people badly loses clients. It does mean that the recruiter's incentives are known and can be computed
(\cref{iq:iv:applications:5}). In the United Kingdom an employment agency may not charge a work-seeker a fee
for finding them work, except in a few occupations listed in regulations (\cref{dat:iv:applications:law});
a recruiter who asks a candidate for money is not a recruiter.

\section{Automated screening}

Large firms receive more applications than people can read, and some screens are partly automated: a program
ranks CVs, or scores a recorded video interview, or filters by answers to form questions. Two consequences
follow. The CV must be readable by a program as well as by a person (plain text, standard section names, no
information only in images). And the rules are changing: automated screening tools are regulated in some
places, with audits and notices to candidates.

\begin{dated}{2026-09}{Rules on automated screening}\label{dat:iv:applications:law}
\textbf{New York City}: Local Law 144 of 2021 prohibits employers and employment agencies from using an
automated employment decision tool unless it has had a bias audit within one year of its use, information
about the audit is publicly available, and the required notices have been given to candidates; the city began
enforcing it on 5 July 2023. \textbf{European Union}: the AI Act (Regulation (EU) 2024/1689), Annex III,
point 4(a), classifies as high-risk ``AI systems intended to be used for the recruitment or selection of
natural persons, in particular to place targeted job advertisements, to analyse and filter job applications,
and to evaluate candidates''; the Digital Omnibus on AI (Regulation (EU) 2026/1744, in force since 27 July
2026) moved the application of the high-risk requirements to such systems from 2 August 2026 to 2 December
2027. \textbf{United Kingdom}: section 6(1)(a) of the Employment Agencies Act 1973 prohibits an employment
agency from charging a fee to a person for finding them employment, with exceptions prescribed by the
Conduct of Employment Agencies and Employment Businesses Regulations 2003 (regulation 26 and Schedule 3,
mostly entertainment and modelling).
\end{dated}

\section{Worked answers}

\begin{example}[A project line that survives the first follow-up]\label{ex:iv:applications:project}
\emph{``Your CV says: `Developed a profitable mean-reversion strategy on ETFs.' Talk me through it.''} The line invites the
question that sinks it: profitable after what? The candidate who has done the arithmetic before the interview says:
``Fifty sector and country ETFs, 2010 to 2020, daily data. Before costs it had a Sharpe ratio of 0.9 at 10\% volatility,
so about 9\% a year. It traded its book about a hundred times a year; at 5 basis points a trade that is 5\% a year, which
leaves 4\% and a Sharpe ratio of 0.4. So: a real effect, most of it eaten by turnover, and my next step was to trade less
often.'' The rewritten line follows from the answer: ``Backtested a daily mean-reversion strategy on 50 ETFs
(2010--2020): Sharpe ratio 0.9 before costs, 0.4 after 5\,bp costs; halving turnover kept 0.6.'' The last clause must be
true and tested before it is written; here the candidate's rerun at half the turnover had a gross return of 8.5\% and costs
of 2.5\%, which is 6\% at 10\% volatility, and the line says exactly that.
\end{example}

\begin{example}[Stating a ranking so that it can be checked]\label{ex:iv:applications:ranking}
\emph{``You write that you were `top 1\% in a forecasting competition'. What exactly was it?''} Rankings invite two
checks: the denominator and the leaderboard. ``Twelfth of 2\,000 teams on the private leaderboard, so the top 0.6\%; on
the public leaderboard I was fifth, and I dropped because my last submissions were tuned to the public split'' is a
complete answer: it gives the number, which board, and the lesson about overfitting to a validation set that
\cref{ch:iv:machine-learning} will test again. ``Top 1\%'' alone reads as rounding in the candidate's favour, and
interviewers ask. \Cref{fig:iv:applications:bestof} shows how far the best of many equally good entries is expected
to sit above its true level.
\end{example}

\begin{omfigure}
\begin{tikzpicture}
\begin{semilogxaxis}[width=9.5cm,height=5.2cm,xlabel={\small number of equally skilled entries, $k$},
  ylabel={\small best score above the truth},xmin=1,xmax=5000,ymin=0,ymax=4.5,
  tick label style={font=\footnotesize},grid=major,grid style={gray!20},table/col sep=comma,
  legend cell align=left,legend pos=south east,legend style={font=\scriptsize}]
\addplot[omDef,thick,mark=*,mark size=1.2pt] table[x=k,y=best]{figdata/interviews/03-applications/bestof.csv};
\addplot[omProp,thick,dashed] table[x=k,y=rule]{figdata/interviews/03-applications/bestof.csv};
\legend{expected best,$\sqrt{2\ln k}$}
\end{semilogxaxis}
\end{tikzpicture}

\omcaption{Expected best of $k$ equally skilled entries whose scores carry independent standard normal noise, in noise
standard deviations above the common true score (solid), and the approximation $\sqrt{2\ln k}$ (dashed), which overstates
it at every size shown. With 400 entries the leader is expected to be 2.97 standard deviations above the truth, which is why
public leaderboards and best-of-many backtests shrink when retested. Data: \texttt{fig\_iv\_bestof.py}.}\label{fig:iv:applications:bestof}
\end{omfigure}

\begin{example}[An unusual path]\label{ex:iv:applications:path}
\emph{``You spent two years as a secondary-school teacher between your degree and your master's. Why?''} The interviewer
is checking that the CV's story is coherent and that the candidate owns it. A scoring answer is short, true and connects
forward: ``I wanted to know whether I could explain mathematics to people who did not want to hear it. I could, mostly,
and I learned to check what someone understood rather than what I had said. I missed working on hard problems myself,
which is why I went back for the master's.'' It avoids apology and does not over-explain; the follow-up, if there is one,
is about the master's.
\end{example}

\section{Question bank}

\begin{interviewq}[$\star$ \iqroles{researcher, mle} \iqfirm{any}]\label{iq:iv:applications:1}
Rewrite this CV line so that a research interviewer would want to ask about it: ``Used Python and scikit-learn
to build machine-learning models for cryptocurrency price prediction.''
\end{interviewq}

\begin{interviewq}[$\star$ \iqroles{trader} \iqfirm{market maker}]\label{iq:iv:applications:2}
``Walk me through your CV in one minute.'' What do you say, in what order, and what do you leave out?
\end{interviewq}

\begin{interviewq}[$\star$ \iqroles{developer} \iqfirm{proprietary firm}]\label{iq:iv:applications:3}
Your CV lists C++, Rust, Python, Java, Go, Haskell and CUDA. The interviewer says: ``Let's do this one in
Haskell.'' You have written two hundred lines of Haskell in your life. What should the CV have said, and what
do you say now?
\end{interviewq}

\begin{interviewq}[$\star\star$ \iqroles{researcher, mle} \iqfirm{systematic fund}]\label{iq:iv:applications:4}
You finished third of 400 teams on a data competition's public leaderboard and fortieth on the private one.
If every team had the same true skill and public scores carried independent noise with standard deviation
0.01, by how much would you expect the best public score to exceed the true score? What do you say when an
interviewer asks about the drop?
\end{interviewq}

\begin{interviewq}[$\star\star$ \iqroles{trader, researcher} \iqfirm{any}]\label{iq:iv:applications:5}
Suppose a recruiter is paid 20\% of your first-year base salary if you accept, and your offer's base is
150\,000. What is the fee? You ask for 160\,000, and the recruiter judges that asking puts a 10\% chance on
the placement failing. How does the recruiter's expected fee change, and what does that tell you about the
advice you will get?
\end{interviewq}

\begin{interviewq}[$\star\star$ \iqroles{developer, mle} \iqfirm{bank}]\label{iq:iv:applications:6}
A large bank tells you that its first screen is partly automated. What can you do about it, what should you
not do, and what rights do candidates have where you live?
\end{interviewq}

\begin{interviewq}[$\star\star\star$ \iqroles{researcher} \iqfirm{systematic fund}]\label{iq:iv:applications:7}
Your CV says a strategy you designed had a backtested Sharpe ratio of 2.4 over five years. In the interview
you admit you tried about sixty variants. What would the best of sixty worthless variants be expected to show,
and how do you defend, or retract, the line?
\end{interviewq}

\begin{interviewq}[$\star\star\star$ \iqroles{trader, researcher, developer} \iqfirm{any}]\label{iq:iv:applications:8}
A former classmate works at the firm you most want and offers to refer you. What do you ask her for, what do
you not ask her for, and how does a referral change what the firm does with your application?
\end{interviewq}

\begin{omsources}
\item S.~V.~Burks, B.~Cowgill, M.~Hoffman and M.~Housman, ``The value of hiring through employee referrals'',
\emph{Quarterly Journal of Economics} 130(2), 2015, 805--839.
\item New York City Local Law 144 of 2021 and the Department of Consumer and Worker Protection's page on
automated employment decision tools.
\item Regulation (EU) 2024/1689 (the AI Act), Annex III; Regulation (EU) 2026/1744 (Digital Omnibus on AI),
Official Journal L, 24 July 2026.
\item Employment Agencies Act 1973, section 6; the Conduct of Employment Agencies and Employment Businesses
Regulations 2003, SI 2003/3319, regulation 26.
\item One Quant Book~4, chapter~12 (multiple testing); One Quant Book~17, chapters 28--29 (entry routes,
education pipelines).
\end{omsources}
